\documentclass[aspectratio=169,11pt]{beamer}

%----------------------------------------------------------------------------------------
%	BEAMER THEME
%----------------------------------------------------------------------------------------

\mode<presentation> {
    \usetheme{CambridgeUS}
    \usecolortheme{default}
}

%----------------------------------------------------------------------------------------
%	PACKAGES
%----------------------------------------------------------------------------------------

\usepackage{graphicx}
\usepackage{booktabs}
\usepackage{amsmath, amssymb}
\usepackage{array}
\usepackage{multirow}
\usepackage{tikz}
\usepackage{threeparttable}
\usepackage{hyperref}


\usepackage{lmodern}
\usepackage{ragged2e}     % per \RaggedRight dentro le colonne
\usepackage{xltabular}    % longtable + tabularx (break + width control)
\usepackage{setspace}     % già lo usi per \singlespacing
% ---- citations: use biblatex (NOT natbib) ----
\usepackage[
  backend=biber,
  style=authoryear,
  natbib=true,
  maxcitenames=2,
  maxbibnames=99
]{biblatex}
\addbibresource{references.bib}
%----------------------------------------------------------------------------------------
%	LAYOUT TWEAKS
%----------------------------------------------------------------------------------------

% frametitle: banda rossa CambridgeUS molto più sottile
\setbeamerfont{frametitle}{size=\normalsize}

\makeatletter
\setbeamertemplate{frametitle}{%
  \nointerlineskip%
  \vskip-4pt% alza la banda verso l'alto (più negativo = più su)
  \begin{beamercolorbox}[wd=\paperwidth,ht=2.3ex,dp=0.8ex,center]{frametitle}%
    \usebeamerfont{frametitle}\insertframetitle
  \end{beamercolorbox}%
  \vskip0pt% quasi nessuno spazio sotto la banda
}
\makeatother


% margini orizzontali leggermente ridotti
\setbeamersize{text margin left=3mm,text margin right=3mm}

% Graphics path (facoltativo)
\graphicspath{{./}{figures/}}

%----------------------------------------------------------------------------------------
%	TITLE PAGE SETUP
%----------------------------------------------------------------------------------------

\title[Fixed Income Arbitrage Alpha]{Exploring Alpha Generation through \\ Fixed Income Arbitrages}
\subtitle{Preliminary Empirical Results — All Three Strategies}

\author{Alessio Ottaviani}
\institute[EDHEC Business School]{
    EDHEC Business School \\
    \medskip
    \textit{Supervisor: Prof. Riccardo Rebonato}
}
\date{\today}

%----------------------------------------------------------------------------------------
%	START DOCUMENT
%----------------------------------------------------------------------------------------

\begin{document}
%------------------------------------------------------------


%----------------------------------------------------------------------------------------
%	TITLE PAGE
%----------------------------------------------------------------------------------------

\begin{frame}
    \titlepage
\end{frame}
\begin{frame}[t]{Construction of Strategy Returns}
\small

\textbf{BTP Italia Basis}\\
Returns are built from opening positions whenever the basis between BTP Italia and the matched nominal BTP becomes sufficiently wide (positive or negative).  
Positions remain open until convergence or maturity.  
Daily trade-level returns are aggregated into a single equal–weighted strategy index, resampled to monthly frequency.

\vspace{0.35cm}

\textbf{CDS--Bond Basis}\\
Universe includes all EUR corporate bonds whose issuer belongs to the iTraxx Main at the trade date.  
We only trade \textit{negative} basis (long bond / long protection), for which implementation is liquid and feasible.  
Trades remain open until the basis reverts or the bond approaches maturity.  
Daily returns across all active trades form the CDS–Bond Basis index.

\vspace{0.35cm}

\textbf{iTraxx Index Skew}\\
We use the four EUR indices (Main, Senior Financials, Subordinated Financials, Crossover).  
At each decision date, only the on–the–run series are considered, each with thresholds calibrated to its volatility.  
The most mispriced index (in absolute basis) is traded.  
Returns are aggregated into a single iTraxx Skew index.

\end{frame}

%------------------------------------------------------------
\begin{frame}[t]{Factor Frameworks for Performance Attribution}
\small

We evaluate our three fixed–income arbitrage strategies under the 
established risk–factor models used in the literature:

\begin{itemize}
    \item \textit{Duarte, Longstaff \& Yu (2007), Review of Financial Studies}.\\
    Fixed–income arbitrage benchmark used to assess exposure to equity, 
    credit and Treasury–based relative–value factors.

    \item \textit{Fung \& Hsieh (2004), Financial Analysts Journal}.\\
    Trend–following and option–based factors widely used to benchmark 
    hedge–fund strategies.

    \item \textit{Brooks, Gould \& Richardson (2020), Journal of Fixed Income}.\\
    Active fixed–income premium decomposition: term, credit, EM, inflation–linked, 
    and volatility premia.
\end{itemize}

\vspace{0.3cm}

\scriptsize
US factors are mapped to their EUR counterparts. 
European equity factors use the Fama–French Europe datasets, converted from USD 
into EUR following the methodology of Glück, Hübel \& Scholz (JPM, 2020). 
Trend–following factors (Fung–Hsieh) are global and directly applicable.

\end{frame}

%------------------------------------------------------------
\begin{frame}{Summary Statistics — Daily Returns (Annualized)}
\centering
\vfill

\scriptsize
\setlength{\tabcolsep}{3pt}

\input{../../results/tables/summary_statistics_daily_beamer.tex}

\vfill
\end{frame}



\begin{frame}[t,shrink=18]{SW vs EW — Basic Metrics (Daily, Annualized)}
\centering
\scriptsize
\setlength{\tabcolsep}{2pt} 
\renewcommand{\arraystretch}{0.95}

\input{../../results/tables/SW_EW_basic_metrics_3strategies.tex}

\end{frame}


\begin{frame}[t,shrink=18]{SW vs EW — MPPM Comparison (Monthly)}
\centering
\scriptsize
\setlength{\tabcolsep}{3pt}

\input{../../results/tables/MPPM_analysis_monthly.tex}

\end{frame}


\begin{frame}[t,shrink=10]{SW vs EW — Volatility-Managed Performance (Monthly)}
\centering
\scriptsize
\setlength{\tabcolsep}{3pt}

\input{../../results/tables/moreira_muir_performance_monthly.tex}

\end{frame}

\begin{frame}[t,shrink=10]{SW vs EW — Volatility-Managed Robustness (Monthly)}
\centering
\scriptsize
\setlength{\tabcolsep}{3pt}

\input{../../results/tables/moreira_muir_robustness_monthly.tex}

\end{frame}


%------------------------------------------------------------
% DUARTE FACTOR MODEL - TABELLA GENERATA AUTOMATICAMENTE
%------------------------------------------------------------
\input{../../results/tables/Duarte_Presentation_Slide_monthly.tex}

%------------------------------------------------------------
% FUNG & HSIEH
%------------------------------------------------------------
\input{../../results/tables/FungHsieh_Presentation_Slide_monthly.tex}

%------------------------------------------------------------
% ACTIVE FI
%------------------------------------------------------------
\input{../../results/tables/ActiveFI_Presentation_Slide_monthly.tex}

%------------------------------------------------------------
% PCA INTRO SLIDE (Implementation Choices)
%------------------------------------------------------------
\begin{frame}[t]{PCA-based Risk Adjustment: Implementation Choices}
\small

\begin{itemize}\itemsep=0.25cm
  \item \textbf{Factor universe:} PCA on the full set of 70+ factors (tradable, non-tradable, macro and transformed variables); no grouping by factor families.
  \item \textbf{Estimation scheme:} out-of-sample rolling PCA with fixed window $W$ (baseline $W=24$ months); loadings estimated on $[t-W,t-1]$.
  \item \textbf{Timing:} both contemporaneous ($R_t$ on $PC_t$) and predictive ($R_{t+1}$ on $PC_t$) specifications.
  \item \textbf{Inference:} Newey--West HAC standard errors; robustness to window length $W$ and number of PCs $K$.
\end{itemize}

\end{frame}


% =========================
% PCA TABLE 1 (Alpha)
% =========================
\begin{frame}[t]{PCA Spanning Regression — Alpha}
\centering
\scriptsize

% ============================================================================
% PCA SPANNING REGRESSION RESULTS
% Rolling Window: 24 months
% Number of PCs: 8
% Timings: predictive, contemporaneous
% ============================================================================

% ============================================================================
% TABLE 1: ALPHA COMPARISON
% ============================================================================

\centering
\begin{threeparttable}
\begin{tabular}{lcccccc}
\toprule
 & \multicolumn{3}{c}{Predictive ($PC_t \to R_{t+1}$)} & \multicolumn{3}{c}{Contemporaneous ($PC_t \to R_t$)} \\
\cmidrule(lr){2-4} \cmidrule(lr){5-7}
Strategy & $\alpha$ (\%) & $R^2$ adj & N & $\alpha$ (\%) & $R^2$ adj & N \\
\midrule
\textit{Btp Italia} & 1.47*** & 0.036 & 157 & 1.37** & 0.040 & 158 \\
 & (2.73) & & & (2.35) & & \\
\addlinespace
\textit{Cds Bond Basis} & 6.38*** & -0.013 & 198 & 6.53*** & 0.073 & 199 \\
 & (4.43) & & & (4.82) & & \\
\addlinespace
\textit{Itraxx Combined} & 3.31*** & 0.002 & 181 & 2.61*** & 0.103 & 182 \\
 & (4.19) & & & (3.92) & & \\
\addlinespace
\bottomrule
\end{tabular}
\begin{tablenotes}
\small
\item \textit{Notes:} Alpha is annualized. PC1--PC8 are the first 8 principal components extracted from 24-month rolling windows. The first 8 PCs explain on average 79.2\% of factor variance (min 73.1\%, max 86.6\%). t-statistics in parentheses based on Newey-West HAC standard errors.
\item *** p$<$0.01, ** p$<$0.05, * p$<$0.10.
\end{tablenotes}
\end{threeparttable}


\end{frame}


% =========================
% PCA TABLE 2 (Full model)
% =========================
\begin{frame}[t,shrink=15]{PCA Spanning Regression — Full Model}

\scriptsize

% ============================================================================
% TABLE 2: FULL MODEL RESULTS (ALPHA + ALL BETAS)
% ============================================================================

\centering
\begin{threeparttable}
\setlength{\tabcolsep}{2pt}
\renewcommand{\arraystretch}{1.05}
\begin{tabular}{lccccccccccc}
\toprule
Strategy & $\alpha$ (\%) & $\beta_{PC1}$ & $\beta_{PC2}$ & $\beta_{PC3}$ & $\beta_{PC4}$ & $\beta_{PC5}$ & $\beta_{PC6}$ & $\beta_{PC7}$ & $\beta_{PC8}$ & $R^2$ adj & N \\
\midrule
\multicolumn{12}{l}{\textbf{Panel A: Predictive ($PC_t \to R_{t+1}$)}} \\
\addlinespace
\textit{Btp Italia} & 1.47*** & 0.006 & 0.011 & -0.022 & -0.026 & -0.058*** & 0.016 & -0.011 & -0.073** & 0.036 & 157 \\
 & (2.73) & (0.45) & (0.78) & (-1.12) & (-0.71) & (-2.80) & (0.55) & (-0.36) & (-2.04) & & \\
\addlinespace
\textit{Itraxx Combined} & 3.31*** & -0.004 & 0.017 & -0.054** & -0.021 & 0.013 & -0.027 & 0.066** & -0.002 & 0.002 & 181 \\
 & (4.19) & (-0.26) & (0.80) & (-2.03) & (-0.76) & (0.48) & (-0.78) & (2.27) & (-0.08) & & \\
\addlinespace
\textit{Cds Bond Basis} & 6.38*** & -0.006 & -0.037 & -0.038 & 0.015 & -0.011 & 0.035 & 0.030 & 0.025 & -0.013 & 198 \\
 & (4.43) & (-0.27) & (-1.55) & (-1.00) & (0.30) & (-0.33) & (1.03) & (0.49) & (0.53) & & \\
\addlinespace
\midrule
\multicolumn{12}{l}{\textbf{Panel B: Contemporaneous ($PC_t \to R_t$)}} \\
\addlinespace
\textit{Btp Italia} & 1.37** & 0.004 & 0.029 & 0.013 & 0.012 & 0.027 & 0.019 & 0.047* & 0.032 & 0.040 & 158 \\
 & (2.35) & (0.39) & (1.62) & (0.47) & (0.46) & (1.03) & (0.82) & (1.75) & (1.29) & & \\
\addlinespace
\textit{Itraxx Combined} & 2.61*** & -0.022** & 0.005 & 0.004 & 0.021 & -0.008 & 0.073*** & 0.072*** & -0.011 & 0.103 & 182 \\
 & (3.92) & (-2.24) & (0.22) & (0.15) & (0.60) & (-0.27) & (2.64) & (2.68) & (-0.53) & & \\
\addlinespace
\textit{Cds Bond Basis} & 6.53*** & -0.024 & -0.043 & 0.059 & -0.099*** & 0.006 & -0.081** & 0.023 & 0.070 & 0.073 & 199 \\
 & (4.82) & (-0.96) & (-1.35) & (1.11) & (-3.04) & (0.10) & (-2.05) & (0.47) & (1.42) & & \\
\addlinespace
\bottomrule
\end{tabular}
\begin{tablenotes}
\tiny
\item This table reports spanning regression results of monthly strategy excess returns on the first 8 principal components. Principal components are extracted from a rolling window of 24 months. Alpha is reported in annualized percentage terms. t-statistics in parentheses are based on Newey-West HAC standard errors. The first 8 PCs explain on average 79.2\% of factor variance (min 73.1\%, max 86.6\%). \item *** p$<$0.01, ** p$<$0.05, * p$<$0.10.
\end{tablenotes}
\vspace{0.05cm}
\tiny
\end{threeparttable}


\end{frame}

\begin{frame}[t]{PCA Scree Plot}
\vspace{-0.3cm}
\centering
\includegraphics[height=0.9\textheight,keepaspectratio]{../../results/figures/pca/pca_scree_plot.png}
\end{frame}

% =========================
% PCA ROBUSTNESS (W, K)
% =========================
\begin{frame}[t,shrink=12]{PCA Robustness — Window Length and Number of PCs}
\centering
\input{../../results/robustness/tables/PCA_robustness_alpha_R2.tex}
\end{frame}



%------------------------------------------------------------
\begin{frame}[t,shrink=15]{MPPM and Volatility-Managed Performance}
\centering
\scriptsize
\setlength{\tabcolsep}{3pt}

% --------- MPPM TABLE ----------
\textbf{Manipulation-Proof Performance Measure (MPPM)}\\[0.1cm]

\begin{threeparttable}
\begin{tabular}{lrrrrrrr}
\toprule
Strategy & \multicolumn{3}{c}{MPPM (\% p.a.)} & Avg Return & Vol & Sharpe & N \\
\cmidrule(lr){2-4}
& $\rho=2$ & $\rho=3$ & $\rho=4$ & (\% p.a.) & (\% p.a.) & & \\
\midrule
CDS\textit{-}Bond Basis   & 4.86 & 4.82 & 4.79 & 4.93 & 2.62 & 1.89 & 207 \\
iTraxx \textit{combined}  & 3.38 & 3.37 & 3.35 & 3.42 & 1.84 & 1.86 & 243 \\
BTP\textit{-}Italia       & 1.58 & 1.57 & 1.55 & 1.62 & 1.84 & 0.88 & 163 \\
\bottomrule
\end{tabular}
\begin{tablenotes}
\tiny
\item MPPM as in Goetzmann et al.\ (2007), based on monthly returns and reported in annualized percentage terms. All moments are annualized.
\end{tablenotes}
\end{threeparttable}

\vspace{0.25cm}

% --------- MOREIRA–MUIR TABLE ----------
\textbf{Volatility-Managed Portfolios (Moreira--Muir, 2017)}\\[0.1cm]

\begin{threeparttable}
\begin{tabular}{lrrrrrrrrrr}
\toprule
Strategy & \multicolumn{3}{c}{Buy-and-Hold} & \multicolumn{3}{c}{Volatility-Managed} & \multicolumn{3}{c}{Regression} & N \\
\cmidrule(lr){2-4} \cmidrule(lr){5-7} \cmidrule(lr){8-10}
& Return & Vol & Sharpe & Return & Vol & Sharpe & Alpha & Beta & $R^2$ & \\
& (\% p.a.) & (\% p.a.) & & (\% p.a.) & (\% p.a.) & & (\% p.a.) & & & \\
\midrule
CDS\textit{-}Bond Basis   & 4.92 & 2.62 & 1.88 & 3.13 & 2.62 & 1.19 & +2.00 & 0.23 & 0.053 & 202 \\
iTraxx \textit{combined}  & 3.43 & 1.84 & 1.86 & 1.63 & 1.84 & 0.89 & -0.04 & 0.49 & 0.239 & 242 \\
BTP\textit{-}Italia       & 1.50 & 1.80 & 0.84 & 1.08 & 1.80 & 0.60 & +0.10 & 0.65 & 0.422 & 162 \\
\bottomrule
\end{tabular}
\begin{tablenotes}
\tiny
\item Volatility-managed portfolios scale exposure inversely with realized variance so that unconditional volatility matches buy-and-hold. Regression alpha measures the volatility-timing premium from $r^{VM}_{t+1} = \alpha + \beta r_{t+1} + \varepsilon_{t+1}$.
\end{tablenotes}
\end{threeparttable}

\vspace{0.15cm}

\tiny
\begin{itemize}
  \item All three strategies deliver positive MPPM, with the highest utility-based performance for the CDS--Bond basis and iTraxx strategies.
  \item Volatility timing generates a sizable premium for the CDS--Bond basis, while effects are more modest for BTP Italia and iTraxx, consistent with their lower return variance.
\end{itemize}

\end{frame}

%------------------------------------------------------------
\begin{frame}[t]{Volatility Timing Robustness — Leverage Constraints}
\centering
\scriptsize
\setlength{\tabcolsep}{4pt}

\begin{threeparttable}
\begin{tabular}{llrrr}
\toprule
Strategy & Metric & No Leverage & 50\% Leverage & Table V \\
         &        & $(\leq 1)$  & $(\leq 1.5)$   & $[0.1,2.0]$ \\
\midrule
CDS\textit{-}Bond Basis   & Alpha (\% p.a.) & +0.60 & +0.80 & +0.83 \\
                          & Sharpe          & 1.53  & 1.48  & 1.49  \\
\addlinespace
iTraxx \textit{combined}  & Alpha (\% p.a.) & +0.03 & +0.01 & -0.02 \\
                          & Sharpe          & 1.35  & 1.18  & 1.09  \\
\addlinespace
BTP\textit{-}Italia       & Alpha (\% p.a.) & +0.02 & +0.14 & +0.21 \\
                          & Sharpe          & 0.69  & 0.73  & 0.74  \\
\bottomrule
\end{tabular}
\begin{tablenotes}
\tiny
\item Leverage constraints follow Moreira and Muir (2017), Table V. ``No Leverage'' constrains the scaling factor to $[0,1]$, ``50\% Leverage'' to $[0,1.5]$, and ``Table V'' to $[0.1,2.0]$. Alpha measures the volatility-timing premium under each constraint.
\end{tablenotes}
\end{threeparttable}

\vspace{0.25cm}

\tiny
\begin{itemize}
  \item Volatility-timing alpha for the CDS--Bond basis remains positive and stable across all leverage constraints, indicating that results are not driven by extreme leverage.
  \item For BTP Italia and iTraxx, volatility timing adds a smaller but consistent premium, with Sharpe ratios largely preserved under realistic constraints.
\end{itemize}

\end{frame}

%------------------------------------------------------------
% BRIDGE SLIDE: OLD RESULTS -> NEW WORK
%------------------------------------------------------------
\begin{frame}[t]{Update Since Last Call: Systematic Extension of Risk Adjustment}
\small


\textbf{What has been done since then (new work)}\\
To strengthen the risk-adjustment side beyond the standard benchmarks, I conducted an extensive review of the recent fixed-income and arbitrage literature and \textbf{reconstructed the full historical time-series dataset} for a comprehensive set of candidate risk controls (70+ factors).  Each series was rebuilt consistently with the original paper definitions (transformations, shocks/innovations, lags, and frequency alignment).

\vspace{0.35cm}

\textbf{Factor universe (organized in 6 families)}\\[-0.05cm]
\begin{itemize}\itemsep=0.12cm
    \item \textit{Credit risk} (e.g., CDS indices, excess bond premium, sovereign spreads)
    \item \textit{Liquidity \& funding stress} (e.g., OIS spreads, repo indicators, noise/illiquidity measures)
    \item \textit{Equity / intermediary proxies} (e.g., BAB, value, momentum)
    \item \textit{Volatility} (e.g., VIX/V2X, MOVE, option-implied measures, trend-following)
    \item \textit{Macro uncertainty \& policy} (e.g., uncertainty indices, EPU, financial conditions)
    \item \textit{Rates \& term structure} (e.g., term, slope, swap spreads)
\end{itemize}

\vspace{0.15cm}
\footnotesize
\textit{Next: full factor list (as a “menu”) and then a parsimonious spanning set will be identified for each strategy.}

\end{frame}

%------------------------------------------------------------

\begin{frame}[t,allowframebreaks]{Risk Factors Universe (Literature-Based, Reconstructed Series)}
\scriptsize
\setlength{\tabcolsep}{3pt}
\renewcommand{\arraystretch}{1.05}
\singlespacing

\begin{xltabular}{\linewidth}{@{}
p{0.15\linewidth}
>{\RaggedRight\arraybackslash}X
p{0.10\linewidth}
@{}}
%\caption*{}
%\label{tab:factors_slides}\\
\toprule
\textbf{Factor ID} & \textbf{Description} & \textbf{Reference} \\
\midrule
\endfirsthead

\multicolumn{3}{c}{{\bfseries Table \thetable\ continued}} \\
\toprule
\textbf{Factor ID} & \textbf{Description} & \textbf{Reference} \\
\midrule
\endhead

\midrule
\multicolumn{3}{@{}p{\textwidth}@{}}{\tiny\textit{Note:} Factors are organized in six families. Each series is reconstructed consistently with the original reference (transformations, shocks/innovations, lags, frequency alignment).} \\
\endfoot

\bottomrule
\endlastfoot
% ========== PANEL A: CREDIT ==========
\midrule
\multicolumn{3}{c}{\textbf{Panel A: Credit Risk Factors}} \\
\midrule
BTP\_BUND & Yield differential between the Italian 10-year BTP and the German Bund. & \citet{fausch2018impact} \\ \addlinespace[0.25em]

CDX\_IG & First difference of the CDS spread on the CDX Investment Grade index. & \citet{klaus2009hedge} \\ \addlinespace[0.25em]

CREDIT\_EU & Same construction as CREDIT\_US, computed using Bloomberg euro-area corporate bond yield indices for BBB and AAA ratings. & \citet{fama1989business} \\ \addlinespace[0.25em]

CREDIT\_US & Bond credit risk factor. Yield spread between Moody's BAA and AAA corporate bond indices. & \citet{fama1989business} \\ \addlinespace[0.25em]

DEF\_EU & Same construction as DEF\_US, computed using long-term European corporate bond total returns relative to long-term German Bund total returns. & \citet{fama1993common} \\ \addlinespace[0.25em]

DEF\_US & Bond default risk factor. Total return spread between a long-term U.S. corporate bond portfolio and a long-term U.S. government bond portfolio. & \citet{fama1993common} \\ \addlinespace[0.25em]

EBP & Excess Bond Premium. Residual component of aggregate corporate bond spreads after netting out expected default risk and bond characteristics via cross-sectional spread regressions; proxies time-varying credit risk compensation tied to intermediary constraints. Used with a one-month lag. & \citet{gilchrist2012credit} \\ \addlinespace[0.25em]

ITRX\_MAIN & First difference of the CDS spread on the iTraxx Europe Main index. & \citet{klaus2009hedge} \\ \addlinespace[0.25em]

MAIN\_5/3\_FLATTENER & First difference of the natural logarithm of the ratio between the 5-year and 3-year iTraxx Europe Main CDS spreads. & Market indicator \\ \addlinespace[0.25em]

PB\_EU\_CDS\_1Y & Average 1-year CDS spread of European prime brokers. & \citet{klaus2009hedge} \\ \addlinespace[0.25em]

PB\_EU\_CDS\_5Y & Average 5-year CDS spread of European prime brokers. & \citet{klaus2009hedge} \\ \addlinespace[0.25em]

PB\_US\_CDS\_1Y & Average 1-year CDS spread of U.S. prime brokers. & \citet{klaus2009hedge} \\ \addlinespace[0.25em]

PB\_US\_CDS\_5Y & Average 5-year CDS spread of U.S. prime brokers. & \citet{klaus2009hedge} \\ \addlinespace[0.25em]

SNR-MAIN & First difference of the 5-year CDS spread difference between the iTraxx Senior Financials index and the iTraxx Europe Main index. & Market indicator \\ \addlinespace[0.25em]

XOVER/MAIN & First difference of the natural logarithm of the ratio between the 5-year iTraxx Crossover and 5-year iTraxx Europe Main CDS spreads. & Market indicator \\ \addlinespace[0.25em]

% ========== PANEL B: LIQUIDITY ==========
\midrule
\multicolumn{3}{c}{\textbf{Panel B: Liquidity Factors}} \\
\midrule
$\Delta$FAILS\_PCT\_TSY & First difference of the ratio of Monthly total notional amount of U.S. Treasury securities fails-to-deliver reported by primary dealers, excluding TIPS, to U.S. Total Debt Outstanding, expressing repo fails as a share of total marketable U.S. Treasury debt. & \citet{fleckenstein2014tips} \\ \addlinespace[0.25em]

EURIBOR\_OIS & Euribor--OIS spread. Difference between 3-month Euribor and the 3-month overnight index swap rate based on EONIA or €STR. & \citet{nyborg2014money} \\ \addlinespace[0.25em]

GC-REPO\_T-BILL & Spread between the 3-month U.S. Treasury general-collateral repo rate and the 3-month U.S. Treasury bill yield. & \citet{bai2019cds} \\ \addlinespace[0.25em]

HKM\_IC & Intermediary capital factor. AR(1) innovation in the market-based capital ratio of primary dealers, scaled by the lagged capital ratio. & \citet{he2017intermediary} \\ \addlinespace[0.25em]

HPW\_NOISE & Treasury noise measure of market-wide illiquidity. Root mean squared deviation between observed U.S. Treasury yields and model-implied yields from a fitted smooth zero-coupon yield curve, aggregated across bonds with maturities between 1 and 10 years. & \citet{hu2013noise} \\ \addlinespace[0.25em]

ILLIQ & Change in the Amihud illiquidity measure. First difference of the ratio of absolute daily return to daily dollar volume, capturing shocks to market illiquidity. & \citet{amihud2015pricing} \\ \addlinespace[0.25em]

LIBOR\_OIS & Libor--OIS spread. Difference between 3-month LIBOR and the 3-month overnight index swap rate. & \citet{nyborg2014money} \\ \addlinespace[0.25em]

LIBOR\_REPO\_SHOCK & Defined as the AR(2) residual of the spread, where the spread is the monthly average of daily 3-month U.S. interbank LIBOR minus the 3-month U.S. Treasury general-collateral repo rate. & \citet{asness2013value} \\ \addlinespace[0.25em]

LIQ\_V & Value-weighted return on a 10--1 portfolio formed by sorting stocks on historical liquidity betas. & \citet{pastor2003liquidity} \\ \addlinespace[0.25em]

LIQNT & Pastor--Stambaugh non-traded liquidity factor. Innovation in aggregate market liquidity constructed from a cross-sectional average of stock-level order-flow-induced return-reversal measures, with predictable components removed. & \citet{pastor2003liquidity} \\ \addlinespace[0.25em]

SILLIQ & Stock illiquidity shock. Innovation in the monthly stock-market illiquidity series, defined as the AR(3) residual of the aggregate Amihud illiquidity measure. & \citet{acharya2013liquidity} \\ \addlinespace[0.25em]

TED\_SHOCK\_EU & Same construction as TED\_SHOCK\_US, computed using the monthly average of daily 3-month Euribor minus the 3-month German government bill rate. & \citet{asness2013value} \\ \addlinespace[0.25em]

TED\_SHOCK\_US & Defined as the AR(2) residual of the TED spread, where the TED spread is the monthly average of daily 3-month U.S. interbank LIBOR minus the 3-month U.S. government bill rate. & \citet{asness2013value} \\ \addlinespace[0.25em]

% ========== PANEL C: EQUITY ==========
\midrule
\multicolumn{3}{c}{\textbf{Panel C: Equity Factors}} \\
\midrule
BAB\_EU & Betting Against Beta (Europe). Same construction as BAB\_US, implemented on the European equity universe. & \citet{frazzini2014betting} \\ \addlinespace[0.25em]

BAB\_US & Betting Against Beta (US). Long low-beta stocks and short high-beta stocks; each leg is rescaled to unit beta, forming a near market-neutral zero-cost factor. & \citet{frazzini2014betting} \\ \addlinespace[0.25em]

CMA\_EU & Investment factor (Conservative Minus Aggressive). Return spread between European conservative and aggressive investment portfolios from independent size-investment sorts. & \citet{fama2017international} \\ \addlinespace[0.25em]

GMOM\_EU & Global momentum factor. Return spread between recent winner and loser equity portfolios based on past 12-month returns excluding the most recent month, expressed in euros. & \citet{asness2013value} \\ \addlinespace[0.25em]

GVAL\_EU & Global value factor. Return spread between high and low book-to-market equity portfolios formed within major developed equity regions, expressed in euros. & \citet{asness2013value} \\ \addlinespace[0.25em]

HML\_EU & Value factor (High Minus Low). Return spread between European high and low book-to-market portfolios from independent size--B/M sorts, approximately size-neutral. & \citet{fama2017international}\\ \addlinespace[0.25em]

MKT\_EU & Market excess return on the European value-weighted equity market portfolio over the 1-month German government bill. Returns are expressed in euros. & \citet{fama2017international} \\ \addlinespace[0.25em]

RMW\_EU & Profitability factor (Robust Minus Weak). Return spread between European robust and weak operating profitability portfolios from independent size-profitability sorts. & \citet{fama2017international}\\ \addlinespace[0.25em]

SMB\_EU & Size factor (Small Minus Big). Return spread between diversified European small- and large-cap portfolios from independent size sorts (five-factor construction). & \citet{fama2017international}\\ \addlinespace[0.25em]

UMD\_EU & Momentum factor (Up Minus Down). Monthly return on a European zero-investment winners-minus-losers momentum portfolio. & \citet{carhart1997persistence} \\ \addlinespace[0.25em]

% ========== PANEL D: VOLATILITY ==========
\midrule
\multicolumn{3}{c}{\textbf{Panel D: Volatility Factors}} \\
\midrule
$\Delta$V2X & Change in V2X. Monthly first difference of the V2X index, the European counterpart of VIX measuring 30-day implied volatility on Euro STOXX 50 options. & \citet{chung2019volatility} \\ \addlinespace[0.25em]

$\Delta$VIX & Change in VIX. Monthly first difference of the VIX index, a model-free measure of 30-day implied volatility on S\&P 500 options. & \citet{chung2019volatility} \\ \addlinespace[0.25em]

ATM\_IV\_CDX & 3-month at-the-money implied volatility on options written on the CDX Investment Grade index. & \citet{novales2019splitting} \\ \addlinespace[0.25em]

ATM\_IV\_ITRX & 3-month at-the-money implied volatility on options written on the iTraxx Europe Main index. & \citet{novales2019splitting} \\ \addlinespace[0.25em]

EP\_SVIX\_1M & Equity premium proxy from the SVIX methodology. Option-implied risk-neutral variance of the market simple return over the next month, computed by integrating out-of-the-money index option prices across strikes. & \citet{martin2017expected} \\ \addlinespace[0.25em]

EP\_SVIX\_3M & Same construction as EP\_SVIX\_1M, using options with a 3-month horizon. & \citet{martin2017expected} \\ \addlinespace[0.25em]

IV\_BUND & 3-month at-the-money implied volatility on the 10-year German Bund futures options, based on 100\% moneyness. & \citet{cremers2021treasury} \\ \addlinespace[0.25em]

IV\_TSY & 3-month at-the-money implied volatility on the 10-year U.S. Treasury futures options, based on 100\% moneyness. & \citet{cremers2021treasury} \\ \addlinespace[0.25em]

MOVE & Option-implied U.S. rates volatility. Yield-curve weighted average of at-the-money implied normal yield volatilities from OTC options on 2-, 5-, 10-, and 30-year constant-maturity interest rate swaps; used lagged by one month. & \citet{bansal2022bond} \\ \addlinespace[0.25em]

MOVE2 & Squared MOVE; used lagged by one month to capture nonlinear exposure to option-implied U.S. interest rate volatility. & \citet{bansal2022bond} \\ \addlinespace[0.25em]

PTFSBD & Return on the bond trend-following factor, constructed from lookback straddles on government bond futures. & \citet{fung2001risk} \\[0.25em]

PTFSCOM & Return on the commodity trend-following factor, constructed from lookback straddles on commodity futures. & \citet{fung2001risk} \\[0.25em]

PTFSFX & Return on the currency trend-following factor, constructed from lookback straddles on major currency futures. & \citet{fung2001risk} \\[0.25em]

PTFSIR & Return on the short-term interest rate trend-following factor, constructed from lookback straddles on 3-month interest rate futures. & \citet{fung2001risk} \\[0.25em]

PTFSSTK & Return on the equity index trend-following factor, constructed from lookback straddles on stock index futures. & \citet{fung2001risk} \\[0.25em]


% ========== PANEL E: MACRO ==========
\midrule
\multicolumn{3}{c}{\textbf{Panel E: Macro Factors}} \\
\midrule
$\Delta$UF & Change in financial uncertainty. First difference of the financial uncertainty index, constructed as $\Delta$UM but using a panel of financial market variables. & \citet{ludvigson2021uncertainty} \\ \addlinespace[0.25em]

$\Delta$UM & Change in macroeconomic uncertainty. First difference of the macro uncertainty index, a model-based measure of expected volatility in the unforecastable component of a large panel of macroeconomic variables. & \citet{ludvigson2021uncertainty} \\ \addlinespace[0.25em]

$\Delta$UR & Change in real uncertainty. First difference of the real-activity uncertainty subindex, constructed from real activity variables. & \citet{ludvigson2021uncertainty} \\ \addlinespace[0.25em]

5Y5YREAL\_YIELD & First difference of the spread between the 5y5y interest rate swap rate and the 5y5y inflation swap rate. & Market indicator \\ \addlinespace[0.25em]

BFCI\_EU & Bloomberg Euro Area Financial Conditions Index. Composite Z-score of euro-area money, bond, and equity indicators, standardized to the 1999 to July 1, 2008 pre-crisis period and expressed in standard deviations from pre-crisis conditions. & \citet{osina2019global} \\ \addlinespace[0.25em]

EPU\_EU & Economic Policy Uncertainty index for Europe. Same construction as EPU\_US, implemented on European newspapers and aggregated across European countries. & \citet{baker2016measuring} \\ \addlinespace[0.25em]

EPU\_US & Economic Policy Uncertainty index for the United States. Newspaper-based index constructed from the normalized frequency of articles that jointly reference the economy, uncertainty, and policy-related terms. & \citet{baker2016measuring} \\ \addlinespace[0.25em]

% ========== PANEL F: RATES ==========
\midrule
\multicolumn{3}{c}{\textbf{Panel F: Interest Rate Factors}} \\
\midrule
$\Delta$SLOPE\_EU & Same construction as $\Delta$SLOPE\_US, computed using the 10-year German government bond yield minus the 3-month German government bill yield. & \citet{ferson1991variation} \\ \addlinespace[0.25em]

$\Delta$SLOPE\_US & First difference of the U.S. yield-curve slope, defined as the 10-year Treasury yield minus the 3-month Treasury bill yield. & \citet{ferson1991variation} \\ \addlinespace[0.25em]

EURUSD\_3M\_IV & First difference of the 3-month at-the-money implied volatility on the EUR/USD exchange rate. & Market indicator \\ \addlinespace[0.25em]

SS2Y & Difference between the 2-year EUR fixed swap par rate and the 2-year German Schatz yield. & \citet{collin2001determinants} \\ \addlinespace[0.25em]

SS5Y & Difference between the 5-year EUR fixed swap par rate and the 5-year German Bobl yield. & \citet{collin2001determinants} \\ \addlinespace[0.25em]

SS10Y & Difference between the 10-year EUR fixed swap par rate and the 10-year German Bund yield. & \citet{collin2001determinants} \\ \addlinespace[0.25em]

TERM\_EU & Same construction as TERM\_US, computed using long-term German Bund total returns in excess of the 1-month German government bill return. & \citet{fama1993common} \\ \addlinespace[0.25em]

TERM\_US & Bond term structure factor. Excess total return on long-term U.S. government bonds over the 1-month U.S. Treasury bill return. & \citet{fama1993common} \\ \addlinespace[0.25em]

YSP\_EU & Yield-curve slope factor. Difference between the 5-year and 1-year German government bond yields. & \citet{koijen2017cross} \\ \addlinespace[0.25em]

YSP\_US & Yield-curve slope factor. Difference between the 5-year and 1-year U.S. Treasury yields. & \citet{koijen2017cross} \\ \addlinespace[0.25em]

\end{xltabular}


\end{frame}

\begin{frame}{From a Large Factor Zoo to a Parsimonious Benchmark}
\small

After reconstructing a comprehensive universe of 70+ risk factors, 
the next step is to identify a \textit{small and interpretable set of controls}
that best spans each arbitrage strategy return series.

\vspace{0.3cm}

\textbf{Objective:}

\begin{itemize}
    \item identify the key factors that statistically span strategy returns
    \item test whether alpha survives after proper risk adjustment
    \item avoid overfitting in a high-dimensional setting
\end{itemize}

\vspace{0.3cm}

Because the candidate factor set is large relative to sample length, we rely on a disciplined \textbf{regularized factor-selection procedure}.
\end{frame}

\begin{frame}{Why Adaptive Elastic Net? \citep{zou2009adaptive}} 
  
\small With more than 70 candidate risk factors, standard OLS is not feasible: 
\begin{itemize} 
  \item multicollinearity across controls 
  \item overfitting in finite samples 
  \item unstable estimates and unreliable alpha inference 
\end{itemize} 

\vspace{0.3cm} 

We model each strategy return as: \[ r_{t} = \alpha + \beta' f_{t} + \varepsilon_{t}, \] where $f_t$ is a high-dimensional vector of candidate risk factors. 
\vspace{0.3cm} 
\textbf{Adaptive Elastic Net} \citep{zou2009adaptive} provides: 
\begin{itemize} 
  \item shrinkage (stability) 
  \item automatic variable selection (sparsity) 
  \item interpretable factor benchmark (few selected controls) 
\end{itemize} 

\end{frame}

\begin{frame}{Elastic Net Regularization}
\small

Elastic Net estimates $\beta$ by solving:

\[
\hat{\beta}^{EN}
=
\arg\min_{\beta}
\left\{
\frac{1}{T}\sum_{t=1}^{T}\bigl(r_t-\alpha-\beta^{\prime}f_t\bigr)^2
+\lambda_2\|\beta\|_2^2
+\lambda_1\|\beta\|_1
\right\}.
\]

\vspace{0.3cm}

\begin{itemize}
    \item $\|\beta\|_1$ selects factors (LASSO effect)
    \item $\|\beta\|_2$ stabilizes correlated predictors (ridge effect)
    \item tuning parameters $(\lambda_1,\lambda_2)$ chosen via a BIC-type criterion.
\end{itemize}

\vspace{0.3cm}

This yields a first sparse benchmark spanning the strategy.

\end{frame}

\begin{frame}{Adaptive Elastic Net (Second-Stage Selection)}
\small

Elastic Net is used as a first step to construct adaptive weights:

\[
w_j = \left|\hat{\beta}^{EN}_j\right|^{-\gamma}
\]


\vspace{0.3cm}

The final Adaptive Elastic Net solves:

\[
\hat{\beta}^{AEN}
=
\arg\min_{\beta}
\Biggl\{
\frac{1}{T}\sum_{t=1}^{T}
\bigl(r_t-\alpha-\beta^{\prime}f_t\bigr)^2
+\lambda_2\|\beta\|_2^2
+\lambda_1\sum_{j=1}^{p} w_j|\beta_j|
\Biggr\}.
\]

\vspace{0.3cm}

\textbf{Interpretation:}

\begin{itemize}
    \item important factors are penalized less
    \item irrelevant factors are shrunk to zero
    \item output: a small set of 4--5 selected controls per strategy
\end{itemize}

\end{frame}

\begin{frame}{Stability Selection: Robust Factor Benchmarks}
\small

A single Elastic Net run may select factors that are sample-specific.

\vspace{0.3cm}

We therefore apply \textbf{stability selection}:

\begin{itemize}
    \item resample the data using moving-block bootstrap (6-month blocks),
    \item repeat the full Adaptive Elastic Net procedure,
    \item keep only factors consistently selected across replications.
\end{itemize}

\vspace{0.3cm}

\textbf{Final benchmark:}

\[
S_i = \{\text{top 4--5 factors with highest selection frequency}\}.
\]

\vspace{0.2cm}

This delivers economically stable risk controls.

\end{frame}

\begin{frame}{Credible Alpha Inference: Selection-Inside Bootstrap}
\small

After selecting a stable factor set $S_i$, we estimate:

\[
r_{i,t}=\alpha_i+\beta_i'f_{S_i,t}+\varepsilon_{i,t}.
\]

\vspace{0.3cm}

\textbf{Key issue:} factor selection creates post-selection bias.

\vspace{0.3cm}

Solution: \textbf{selection-inside bootstrap}

\begin{itemize}
    \item resample returns and factors in blocks,
    \item re-run factor selection inside each draw,
    \item obtain confidence intervals for $\hat{\alpha}_i$.
\end{itemize}

\vspace{0.2cm}

This ensures alpha significance is robust to model uncertainty.

\end{frame}


%------------------------------------------------------------
% AEN STABLE OLS - AUTO GENERATED SLIDES
%------------------------------------------------------------
\input{../../results/aen/hqc/tables/AEN_Stable_OLS_BTP_Italia_Presentation_Slide.tex}
\input{../../results/aen/hqc/tables/AEN_Stable_OLS_CDS_Bond_Basis_Presentation_Slide.tex}
\input{../../results/aen/hqc/tables/AEN_Stable_OLS_iTraxx_Combined_Presentation_Slide.tex}
\end{document}
